% 기어식 제어 보드의 시간 신호 경로 (그림 fig:hw-gear-gate). hardware/gear/control-board/README.md,
% 부록 D의 기어식 제어 보드 절. STROBE = S ∧ T, CLEAR = C ∧ T, JUDGE = J ∧ SEEN ∧ T, T = TXEN ∨ CYC.
\begin{tikzpicture}[x=0.96cm,y=1cm,
    b/.style={blk,minimum width=17mm,minimum height=7mm,font=\scriptsize},
    g/.style={blk,minimum width=8mm,minimum height=7mm,font=\small,fill=cbluebg,draw=cblue},
    ar/.style={->,line width=0.5pt}]
  \foreach \y/\s/\n/\o in {2.4/J/JUDGE/34, 1.2/S/STROBE/1, 0/C/CLEAR/35} {
    \node[font=\footnotesize,anchor=east] (in\s) at (0,\y) {센서 $\s$};
    \node[b] (f\s) at (1.3,\y) {RC 필터\\슈미트};
    \node[g] (a\s) at (3.6,\y) {$\wedge$};
    \node[b,fill=corangebg,draw=corange] (d\s) at (5.8,\y) {구동기\\자동 복귀 퓨즈};
    \node[font=\footnotesize,anchor=west] (o\s) at (7.2,\y) {\code{\n} (제어 커넥터 \o번)};
    \draw[ar] (in\s) -- (f\s);
    \draw[ar] (f\s) -- (a\s);
    \draw[ar] (a\s) -- (d\s);
    \draw[ar] (d\s) -- (o\s);
  }
  % T = TXEN ∨ CYC
  \node[g] (orT) at (2.0,-1.5) {$\vee$};
  \node[font=\scriptsize,anchor=east] (tx) at (1.2,-1.25) {\textit{TXEN} (Pico의 허가)};
  \node[font=\scriptsize,anchor=east] (cy) at (1.2,-1.75) {\textit{CYC} (시작한 주기는 끝까지, L7)};
  \draw[ar] (tx.east) -- (orT.west |- tx.east);
  \draw[ar] (cy.east) -- (orT.west |- cy.east);
  \draw (orT.east) -- (3.0,-1.5) coordinate (tb);
  \node[font=\scriptsize,anchor=north] at (2.7,-1.55) {$T$};
  \draw[ar] (tb) |- ([yshift=-1.5mm]aC.west);
  \draw[ar] (tb) |- ([yshift=-1.5mm]aS.west);
  \draw[ar] (tb) |- ([yshift=-1.5mm]aJ.west);
  % SEEN
  \node[b,fill=cgreenbg,draw=cgreen] (seen) at (3.6,3.6) {\textit{SEEN}: 이번 주기에\\\code{STROBE}가 있었다};
  \draw[ar] (seen.south) -- (aJ.north);
  \draw[ar] (4.6,1.2) |- (seen.east);
  \fill (4.6,1.2) circle (1.2pt);
  \node[font=\scriptsize,text=cgray,anchor=west,align=left] at (4.9,3.6) {지우는 때: 전원 투입, \code{HOME},\\정지 래치가 풀릴 때 (L8)};
  \node[note,anchor=north west,align=left] at (3.4,-1.1) {펌웨어는 ``언제 허가할지''만 정하고 펄스의 모양은 슬롯이 정한다.\\
    \code{CLEAR}에는 \textit{SEEN}이 없으므로 \code{HOME} 끝에 \code{CLEAR}만 한 번 낼 수 있다.};
\end{tikzpicture}
